\section{Conclusiones y Recomendaciones}
\label{sec:conclusiones}

\subsection{Balance de objetivos}
\begin{enumerate}
    \item \textbf{OE1. Modelado.} Se formuló el grafo vial dirigido, la función de costo temporal con pendiente y el factor climático uniforme. La no negatividad de los pesos quedó garantizada por construcción, y la instancia ilustrativa se separó explícitamente del levantamiento geográfico.
    \item \textbf{OE2. Implementación.} Se implementaron desde cero el heap binario indexado, el heap 4-ario indexado y el Fibonacci heap, con un contrato común que permite sustituirlos dentro de Dijkstra sin modificar el solver. También se implementaron $A^*$ con heurística consistente, la referencia de Bellman-Ford y el certificado de optimalidad.
    \item \textbf{OE3. Validación.} La suite automatizada reúne 246 pruebas y finalizó con \texttt{246 passed}. Las tres colas produjeron las mismas distancias en todas las instancias, el certificado de Bellman se cumplió y la comparación con NetworkX no detectó errores.
    \item \textbf{OE4. Evaluación.} Se ejecutó un benchmark canónico de 1350 muestras cronometradas, una consulta sobre una red OSM real de 891 vértices y 2139 aristas, una extensión $A^*$ y un análisis de sensibilidad climática.
\end{enumerate}

\subsection{Respuesta a las preguntas de investigación}
\paragraph{PI1. ¿Qué cola minimiza el tiempo y cómo cambia con la escala?} No hay una cola ganadora en todas las escalas. El heap 4-ario obtuvo la menor mediana en 7 de 10 escalas sintéticas, con una razón mediana de 0.964 en el perfil disperso y 0.954 en el denso. El heap binario ganó 3 escalas y la consulta OSM. Los IQR se solaparon en las 10 escalas sintéticas, de modo que la ventaja del 4-ario es pequeña y dependiente del tamaño. En la consulta OSM, el binario fue 26.9\,\% más rápido que el 4-ario y 66.7\,\% más rápido que el Fibonacci.

\paragraph{PI2. ¿La ventaja asintótica del Fibonacci heap se traduce en tiempo de pared?} No en esta implementación. El Fibonacci heap no ganó ninguna escala y fue entre 1.05 y 1.59 veces más lento que el binario. Realizó menos comparaciones (19.5\,\% menos en $|V|=5000$ disperso y 6.7\,\% menos en $|V|=800$ denso), pero el costo de asignación, punteros, cortes y consolidación dominó el tiempo. La cota $O(|E|+|V|\log|V|)$ sigue siendo válida como análisis asintótico; no es una predicción de rendimiento en Python.

\paragraph{PI3. ¿Se reproduce el comportamiento en la red real y qué aportan $A^*$ y el clima?} La consulta OSM confirmó que la elección depende del tamaño: en una red pequeña, el binario fue el más rápido. $A^*$ redujo las expansiones de 887 a 823 (7.2\,\%) sin cambiar el costo. El clima incrementó el tiempo estimado un 6.91\,\%, pero no cambió la ruta, porque un factor uniforme multiplica todos los pesos por la misma constante.

\subsection{Conclusión principal}
La conclusión central es que \emph{la elección de cola es una decisión empírica y condicionada}. La teoría asintótica ordena las estructuras para el peor caso, pero no determina cuál es más rápida en una red vial concreta, en un lenguaje concreto y en una máquina concreta. En este estudio, el heap binario es la opción por defecto más robusta; el heap 4-ario es un candidato con ventaja pequeña en escalas medianas y grandes; el Fibonacci heap no es competitivo en Python. Esa conclusión no se apoya en una sola escala ni en una sola consulta, sino en el patrón conjunto de 10 escalas sintéticas y una consulta real.

\subsection{Contribuciones}
\begin{itemize}
    \item Un núcleo algorítmico modular con tres colas intercambiables, $A^*$ y referencias independientes.
    \item Un protocolo de medición reproducible con semillas, calentamiento, GC desactivado, mediana e IQR, y datos versionados.
    \item Un caso aplicado sobre OpenStreetMap con consulta determinista, ruta reconstruida y certificado de optimalidad.
    \item Una aplicación Streamlit integrada que permite comparar resultados, explorar la ruta, revisar la traza manual y consultar la procedencia de los datos.
    \item Un servicio meteorológico con \textit{snapshot} y \textit{fallback} sin red, separado del benchmark principal.
    \item Una suite de 246 pruebas que cubre contrato de colas, optimalidad, entradas inválidas, CLI, OSM, clima y aplicación.
\end{itemize}

\subsection{Recomendaciones}
\paragraph{Para la práctica.}
\begin{enumerate}
    \item Usar el heap binario como valor por defecto en una aplicación nueva, salvo que un benchmark local indique lo contrario.
    \item Probar el heap 4-ario con la distribución real de consultas y exigir que la ventaja sea mayor que la variabilidad (IQR) antes de adoptarlo.
    \item No adoptar el Fibonacci heap en Python por su cota teórica; reservarlo para análisis de peor caso o para entornos con asignación muy barata.
    \item Usar $A^*$ cuando haya coordenadas y destino conocido; verificar la consistencia de la heurística antes de confiar en la optimalidad.
\end{enumerate}

\paragraph{Para el trabajo futuro.}
\begin{enumerate}
    \item Medir memoria residente y asignaciones para completar la comparación espacial.
    \item Repetir el benchmark en C++ o Rust y en otra CPU para evaluar la portabilidad de las conclusiones.
    \item Añadir pruebas de significancia con intervalos de confianza o \textit{bootstrap}, y aumentar el número de instancias por escala.
    \item Modelar el clima por tramo, con pendiente, tipo de vía y exposición, para que el factor pueda alterar la ruta.
    \item Extender la consulta OSM a un conjunto de pares origen-destino y a distintas horas, para medir la variabilidad entre consultas.
    \item Incorporar congestión dinámica y ventanas temporales, manteniendo la no negatividad de los pesos.
\end{enumerate}

\subsection{Cierre}
El proyecto cumple los objetivos dentro de su alcance y deja explícitas sus limitaciones. La aportación más útil no es un ranking definitivo, sino un método para decidir con evidencia: implementar las alternativas, verificar la corrección, medir con dispersión y repetir la medición en el entorno de despliegue. Esa es la lección transferible a otros problemas de caminos mínimos y a otras estructuras de datos.
